\documentclass{cv}

\hypersetup{
    pdftitle={Jeongjun Ahn's CV},
    pdfauthor={Jeongjun Ahn}
}
\cvname{Jeongjun Ahn}
\cvlastupdated{Last updated in July 2026}

\begin{document}

    \placelastupdatedtext
    \begin{header}
        \textbf{\fontsize{24 pt}{24 pt}\selectfont Jeongjun Ahn}

        \vspace{0.3 cm}

        \normalsize
        \mbox{{\color{black}\footnotesize\faMapMarker*}\hspace*{0.13cm}Leipzig, Germany}%
        \kern 0.25 cm%
        \AND%
        \kern 0.25 cm%
        \mbox{\hrefWithoutArrow{mailto:tyyyu1268@gmail.com}{\color{black}{\footnotesize\faEnvelope[regular]}\hspace*{0.13cm}tyyyu1268@gmail.com}}%
        \kern 0.25 cm%
        \AND%
        \kern 0.25 cm%
        \mbox{\hrefWithoutArrow{tel:+49 152 31627865}{\color{black}{\footnotesize\faPhone*}\hspace*{0.13cm}0152 31627865}}%
        \kern 0.25 cm%
        \AND%
        \kern 0.25 cm%
        \mbox{\hrefWithoutArrow{https://linkedin.com/in/jeong-jun-ahn}{\color{black}{\footnotesize\faLinkedinIn}\hspace*{0.13cm}jeong-jun-ahn}}%
        \kern 0.25 cm%
        \AND%
        \kern 0.25 cm%
        \mbox{\hrefWithoutArrow{https://github.com/prosk-sudo}{\color{black}{\footnotesize\faGithub}\hspace*{0.13cm}prosk-sudo}}%
    \end{header}

    \vspace{0.3 cm - 0.3 cm}


    \section{Education}

        \begin{twocolentry}{
        \textit{Oct 2023 – May 2026}}
            \textbf{Lancaster University Leipzig}

            \textit{BSc in Computer Science}
        \end{twocolentry}

    \section{Languages}

        \begin{onecolentry}
            \textbf{Korean} (Native) \textperiodcentered\ \textbf{English} (Fluent) \textperiodcentered\ \textbf{German} (Beginner)
        \end{onecolentry}

    \section{Projects}

        \begin{twocolentry}{
        \textit{Mar 2026}}
            \textbf{DOMtutor - DOMjudge v9 Migration}

            \textit{BSc Dissertation \textperiodcentered\ \href{https://github.com/DOMtutor}{github.com/DOMtutor}}
        \end{twocolentry}

        \vspace{0.10 cm}
        \begin{onecolentry}
            \begin{highlights}
                \item Migrated DOMtutor, a suite of automation tools for the open-source DOMjudge programming-contest judge, from DOMjudge v7.3.5 to v9.0.0 across 16 releases of breaking changes, following a fork-based incremental strategy.
                \item Rewrote the pyjudge library's raw SQL queries and data models to match a substantially reorganised database schema - restructured executable storage, renamed and removed columns, new required fields and external identifiers, and a new task-based judging pipeline.
                \item Built a Dockerised environment running DOMjudge v7 and v9 side by side for direct comparison and debugging, and updated the Kattis-format problemtools integration for problem verification and PDF generation.
                \item Verified all core workflows - settings, problem, and contest uploads, and submission judging - end-to-end against a live DOMjudge v9 instance.
            \end{highlights}
        \end{onecolentry}

        \vspace{0.13 cm}

        \begin{twocolentry}{
        \textit{2025}}
            \textbf{Lancaster Under Attack}

            \textit{\href{https://github.com/prosk-sudo/Lancaster-Under-Attack}{Lancaster-Under-Attack}}
        \end{twocolentry}

        \vspace{0.10 cm}
        \begin{onecolentry}
            \begin{highlights}
                \item Developed a top-down zombie-survival shooter game in Java using libGDX framework, as a second-year group project, with multiple levels, mouse-aimed shooting, and objectives to eliminate enemies and rescue NPCs.
            \end{highlights}
        \end{onecolentry}

        \vspace{0.13 cm}

        \begin{twocolentry}{
        \textit{2026}}
            \textbf{codewars.nvim}

            \textit{\href{https://github.com/prosk-sudo/codewars.nvim}{codewars.nvim}}
        \end{twocolentry}

        \vspace{0.10 cm}
        \begin{onecolentry}
            \begin{highlights}
                \item Created a Neovim plugin in Lua that lets users browse and solve Codewars katas from codewars.com without leaving the editor, inspired by the plugin leetcode.nvim.
            \end{highlights}
        \end{onecolentry}

        \vspace{0.13 cm}

        \begin{twocolentry}{
        \textit{2026}}
            \textbf{linkedin-cli}

            \textit{\href{https://github.com/prosk-sudo/linkedin-cli}{linkedin-cli}}
        \end{twocolentry}

        \vspace{0.10 cm}
        \begin{onecolentry}
            \begin{highlights}
                \item Revived an abandoned open-source Python command-line client for the LinkedIn API v2 (originally by tigillo), migrating its authentication to the required OpenID Connect.
            \end{highlights}
        \end{onecolentry}

        \vspace{0.13 cm}

        \begin{twocolentry}{
        \textit{2026}}
            \textbf{music-scores}

            \textit{\href{https://github.com/prosk-sudo/music-scores}{music-scores}}
        \end{twocolentry}

        \vspace{0.10 cm}
        \begin{onecolentry}
            \begin{highlights}
                \item Transcribe and typeset sheet music using LilyPond as an ongoing personal project.
            \end{highlights}
        \end{onecolentry}

    \section{Leadership \& Activities}

        \begin{twocolentry}{
        \textit{Nov 2023}}
            \textbf{IBM Qiskit Fall Fest}
        \end{twocolentry}

        \vspace{0.10 cm}
        \begin{onecolentry}
            \begin{highlights}
                \item Took part in IBM's global quantum computing hackathon series, working with Qiskit to explore quantum computing concepts.
            \end{highlights}
        \end{onecolentry}

        \vspace{0.13 cm}

        \begin{twocolentry}{
        \textit{Jan 2024}}
            \textbf{Global Game Jam 2024}

            \textit{\href{https://github.com/LUComp/ggj2024}{LUComp/ggj2024}}
        \end{twocolentry}

        \vspace{0.10 cm}
        \begin{onecolentry}
            \begin{highlights}
                \item Collaborated with a group of people from the university to design, build, and playtest a game within the event's 48-hour format, with the theme "Make Me Laugh".
            \end{highlights}
        \end{onecolentry}

        \vspace{0.13 cm}

        \begin{twocolentry}{
        \textit{Jan 2025}}
            \textbf{Global Game Jam 2025}

            \textit{\href{https://github.com/LUComp/ggj2025}{LUComp/ggj2025}}
        \end{twocolentry}

        \vspace{0.10 cm}
        \begin{onecolentry}
            \begin{highlights}
                \item Collaborated with a group of people from the university to design, build, and playtest a game within the event's 48-hour format, with the theme "Bubbles".
            \end{highlights}
        \end{onecolentry}

        \vspace{0.13 cm}

        \begin{twocolentry}{
        \textit{Oct 2024 – May 2026}}
            \textbf{LUComp \& Robotics Club} - \textit{Member}
        \end{twocolentry}

        \vspace{0.10 cm}
        \begin{onecolentry}
            \begin{highlights}
                \item Regularly attended club meetings and events, engaging with fellow members on computing and robotics projects.
            \end{highlights}
        \end{onecolentry}

        \vspace{0.13 cm}

        \begin{twocolentry}{
        \textit{Nov 2024 – May 2026}}
            \textbf{Chess Club, Lancaster University Leipzig} - \textit{President}
        \end{twocolentry}

        \vspace{0.10 cm}
        \begin{onecolentry}
            \begin{highlights}
                \item Organised and led the university's chess club as President, coordinating regular sessions and matches for members.
            \end{highlights}
        \end{onecolentry}

        \vspace{0.13 cm}

        \begin{twocolentry}{
        \textit{Oct 2025 – May 2026}}
            \textbf{Deputy Student Representative}
        \end{twocolentry}

        \vspace{0.10 cm}
        \begin{onecolentry}
            \begin{highlights}
                \item Served as a liaison between students and the university, helping to raise and resolve academic and administrative concerns.
            \end{highlights}
        \end{onecolentry}

    \section{Technologies}

        \begin{onecolentry}
            \textbf{Programming Languages:} Python, C, Java, JavaScript, Lua, SQL, HTML/CSS, Lilypond
        \end{onecolentry}

        \vspace{0.13 cm}

        \begin{onecolentry}
            \textbf{Frameworks \& Libraries:} FastAPI, PyTorch, scikit-learn, Hugging Face, LibGDX
        \end{onecolentry}

        \vspace{0.13 cm}

        \begin{onecolentry}
            \textbf{Tools \& Platforms:} Linux, Bash, Git, GitHub, Docker, uv
        \end{onecolentry}

\end{document}
